\documentclass[11pt]{article}
\usepackage[margin=1in]{geometry}
\usepackage{amsmath}
\usepackage{amssymb}

\begin{document}

\begin{center}
\textbf{\Large Homework 6} \\
Tashfeen Hatim Said Omran \\
UCF ID: 5862894 \\
Class: STA4321 \\
October 22, 2026
\end{center}

\vspace{0.5cm}

\noindent \textbf{Problem 4.17} \\
a) I find $c$ by integrating the density function over its domain and setting it equal to 1:
\[ \int_{0}^{1} (cy^2 + y) \, dy = 1 \]
\[ \left[ c\frac{y^3}{3} + \frac{y^2}{2} \right]_0^1 = 1 \implies \frac{c}{3} + \frac{1}{2} = 1 \implies \frac{c}{3} = \frac{1}{2} \implies c = \frac{3}{2} \]

\noindent b) I find $F(y)$ by integrating the density function $f(t)$ from $0$ to $y$:
\[ F(y) = \int_{0}^{y} \left( \frac{3}{2}t^2 + t \right) \, dt = \frac{1}{2}y^3 + \frac{1}{2}y^2 \]
For $y < 0$, $F(y) = 0$. For $y > 1$, $F(y) = 1$.

\noindent c) I plot $f(y) = 1.5y^2 + y$ as a parabola opening upwards from the origin to the point $(1, 2.5)$. I plot $F(y) = 0.5y^3 + 0.5y^2$ as an increasing curve from $(0,0)$ to $(1,1)$. 

\noindent d) I use $F(y)$ to find the values:
\[ F(-1) = 0 \]
\[ F(0) = 0 \]
\[ F(1) = \frac{1}{2}(1)^3 + \frac{1}{2}(1)^2 = 1 \]

\noindent e) I find the probability of finishing in less than half an hour ($Y < 0.5$):
\[ P(Y < 0.5) = F(0.5) = \frac{1}{2}(0.5)^3 + \frac{1}{2}(0.5)^2 = \frac{1}{16} + \frac{1}{8} = \frac{3}{16} = 0.1875 \]

\noindent f) I find the conditional probability $P(Y \ge 0.5 \mid Y \ge 0.25)$:
\[ P(Y \ge 0.5 \mid Y \ge 0.25) = \frac{P(Y \ge 0.5)}{P(Y \ge 0.25)} = \frac{1 - F(0.5)}{1 - F(0.25)} \]
\[ F(0.25) = \frac{1}{2}(0.25)^3 + \frac{1}{2}(0.25)^2 = \frac{1}{128} + \frac{1}{32} = \frac{5}{128} \]
\[ P(Y \ge 0.5 \mid Y \ge 0.25) = \frac{1 - \frac{3}{16}}{1 - \frac{5}{128}} = \frac{\frac{13}{16}}{\frac{123}{128}} = \frac{104}{123} \approx 0.8455 \]

\vspace{0.5cm}
\noindent \textbf{Problem 4.21} \\
I find the mean, $E(Y)$:
\[ E(Y) = \int_{0}^{1} y \left( \frac{3}{2}y^2 + y \right) \, dy = \int_{0}^{1} \left( \frac{3}{2}y^3 + y^2 \right) \, dy = \left[ \frac{3}{8}y^4 + \frac{1}{3}y^3 \right]_0^1 = \frac{3}{8} + \frac{1}{3} = \frac{17}{24} \]

\noindent I find the variance, $V(Y) = E(Y^2) - [E(Y)]^2$:
\[ E(Y^2) = \int_{0}^{1} y^2 \left( \frac{3}{2}y^2 + y \right) \, dy = \int_{0}^{1} \left( \frac{3}{2}y^4 + y^3 \right) \, dy = \left[ \frac{3}{10}y^5 + \frac{1}{4}y^4 \right]_0^1 = \frac{3}{10} + \frac{1}{4} = \frac{11}{20} \]
\[ V(Y) = \frac{11}{20} - \left(\frac{17}{24}\right)^2 = \frac{11}{20} - \frac{289}{576} = \frac{3168 - 2890}{5760} = \frac{278}{5760} = \frac{139}{2880} \]

\vspace{0.5cm}
\noindent \textbf{Problem 4.26} \\
a) I use the definition of expected value:
\[ E(aY + b) = \int_{-\infty}^{\infty} (ay + b)f(y) \, dy = a\int_{-\infty}^{\infty} y f(y) \, dy + b\int_{-\infty}^{\infty} f(y) \, dy = aE(Y) + b = a\mu + b \]

\noindent b) I use the definition of variance, $V(X) = E[(X - E(X))^2]$, and substitute $X = aY + b$:
\[ V(aY + b) = E[((aY + b) - (a\mu + b))^2] = E[(aY - a\mu)^2] \]
\[ V(aY + b) = E[a^2(Y - \mu)^2] = a^2 E[(Y - \mu)^2] = a^2 V(Y) = a^2 \sigma^2 \]

\vspace{0.5cm}
\noindent \textbf{Problem 4.27} \\
I use the results from Exercise 4.21 where $E(Y) = 17/24$ and $V(Y) = 139/2880$, and the theorems from 4.26.
\[ E(W) = E(5 - 0.5Y) = 5 - 0.5E(Y) = 5 - 0.5\left(\frac{17}{24}\right) = 5 - \frac{17}{48} = \frac{223}{48} \]
\[ V(W) = V(5 - 0.5Y) = (-0.5)^2 V(Y) = 0.25\left(\frac{139}{2880}\right) = \frac{139}{11520} \]

\vspace{0.5cm}
\noindent \textbf{Problem 4.38} \\
a) I find $F(y)$ for a uniform distribution over $(0, 1)$:
\[ F(y) = \int_{0}^{y} 1 \, dt = y \quad \text{for } 0 \le y \le 1 \]
$F(y) = 0$ for $y < 0$, and $F(y) = 1$ for $y > 1$.

\noindent b) I show the probability depends only on $b$:
\[ P(a \le Y \le a + b) = F(a + b) - F(a) = (a + b) - a = b \]
This result relies only upon the value of $b$.

\vspace{0.5cm}
\noindent \textbf{Problem 4.50} \\
I define $Y$ as a uniform distribution over the interval $(0, 500)$.
\\
a) I find the probability that it is within 25 feet of the end ($Y > 475$):
\[ P(Y > 475) = \frac{500 - 475}{500} = \frac{25}{500} = 0.05 \]

\noindent b) I find the probability that it is within 25 feet of the beginning ($Y < 25$):
\[ P(Y < 25) = \frac{25 - 0}{500} = \frac{25}{500} = 0.05 \]

\noindent c) I find the probability that it is closer to the beginning than the end ($Y < 250$):
\[ P(Y < 250) = \frac{250 - 0}{500} = 0.50 \]

\vspace{0.5cm}
\noindent \textbf{Problem 4.73} \\
I have a normal distribution with $\mu = 950$ and $\sigma = 10$.
\\
a) I find the probability that the width is between $947$ and $958$:
\[ P(947 < Y < 958) = P\left(\frac{947 - 950}{10} < Z < \frac{958 - 950}{10}\right) = P(-0.3 < Z < 0.8) \]
\[ P(-0.3 < Z < 0.8) = P(Z < 0.8) - P(Z < -0.3) = 0.7881 - 0.3821 = 0.4060 \]

\noindent b) I find $C$ such that $P(Y < C) = 0.8531$:
\[ P\left(Z < \frac{C - 950}{10}\right) = 0.8531 \]
Using the standard normal table, the z-score for $0.8531$ is $1.05$.
\[ \frac{C - 950}{10} = 1.05 \implies C - 950 = 10.5 \implies C = 960.5 \text{ mm} \]

\end{document}